% =====================================================================
%  1 -- THE STAR AS A THERMODYNAMIC ENGINE, AND ITS LAST DAY.
%
%  Source: tier3.md Part 0 header + intro incl. the organising claim
%          (228-254), 0.1 (256-441) with 0.1.1 (262), 0.1.2 (296),
%          0.1.3 (351), 0.1.4 (398); 0.2 (443-869) with 0.2.1 (450),
%          0.2.2 (574), 0.2.3 (612), 0.2.4 (711), 0.2.5 (734).
%          Self-check: source 0.10 (2714-2755) questions Q1, Q2, Q3,
%          Q4, Q5, Q11, Q12 (the rest of that block is redistributed
%          to sections 2, 3, 4 and 6 -- see Appendix B).
%
%  Re-homings / rewordings performed here:
%    - The organising-claim paragraph "Every clause of that has to be
%      earned. SS0.1-0.4 earn ..." keeps its wording; each SS-range is
%      replaced by the report's live \ref ranges (per the brief).
%    - Source typo `\m\={u}` (backslash-m + u-macron) in the T_c equation (line 357/359) -> \bar\mu.
%    - The em-dashes inside \text{} in the mixing equation (754-757)
%      -> ---.
%    - No sentence dropped, merged, or reworded otherwise.
% =====================================================================

\section{The star as a thermodynamic engine, and its last day}
\label{part:engine}

Three earlier documents touch this material and none of them is a substitute for
it. Tier~0 Part~I lists the motivation chain as claims to defend. Tier~2 Part~0
derives the equilibrium hierarchy and the $\Ye \to \Mch$ link because the
algebra of that tier is unreadable without them. This part is the
\emph{setting}: the star, the last day, the collapse, and the light that comes
out. It is written so that someone who has never opened a stellar-structure
textbook can follow every step, and so that someone who has can check every one.

The organising claim, stated once so that everything below can be measured
against it:

\begin{invariant}
\textbf{The project exists because a one-day episode of nuclear burning in a
sphere a few thousand km in radius sets a single scalar --- \Ye{} --- that
propagates through a collapse, a bounce, a shock and an explosion to become a
number an astronomer measures. The episode is stiff, the network is expensive,
and every stellar-evolution calculation currently pays for it by using a
network too small to get \Ye{} right.}
\end{invariant}

% [edition] reworded cross-ref: "SS0.1-0.4 earn ... SS0.5 earns ... SS0.6 earns
% ... SS0.7 earns ... SS0.8 earns" -- the section ranges are replaced by the
% report's live references.
Every clause of that has to be earned. \S\ref{part:engine},
\S\ref{part:siburn} and \S\ref{part:collapse} earn ``one-day episode of nuclear
burning''. \S\ref{part:collapse} earns ``propagates through a collapse''.
\S\ref{part:observables} earns ``a number an astronomer measures''.
\S\ref{part:weak} earns ``sets a single scalar''. \S\ref{part:host} earns
``expensive'' and ``too small''.

% =====================================================================
\subsection{The star as a self-gravitating thermodynamic engine}
\label{sec:engine}

Everything about late stellar evolution follows from one structural fact: a star
is a bound system whose gravity is its own, so its thermodynamics has a sign
that laboratory thermodynamics does not.

% ---------------------------------------------------------------------
\subsubsection{Hydrostatic equilibrium, and the pressure a star needs}
\label{sec:hydrostatic}

A spherically symmetric star in mechanical balance satisfies (textbook:
\citealt{kippenhahn2012}, Ch.~2; \citealt{clayton1968}, Ch.~2)
%
\begin{equation}
\frac{dP}{dr} \;=\; -\,\frac{G\,m(r)\,\rho(r)}{r^{2}},
\qquad
\frac{dm}{dr} \;=\; 4\pi r^{2}\rho(r).
\label{eq:hydrostatic}
\end{equation}

Integrating the first crudely --- replace $dP/dr$ by $-P_c/R$, $m$ by $M$,
$\rho$ by the mean $M/(4\pi R^3/3)$ --- gives the central pressure a star of
mass $M$ and radius $R$ must supply:
%
\begin{equation}
\begin{gathered}
P_c \;\sim\; \frac{3}{4\pi}\,\frac{GM^{2}}{R^{4}}\\[2pt]
\text{(and }\tfrac{3}{8\pi}\text{ if the uniform-density case is
integrated exactly rather than estimated).}
\end{gathered}
\label{eq:centralpressure}
\end{equation}

The exact coefficient depends on the density profile --- the two above differ by
a factor of two and neither is right for a real star --- so the scaling
$P_c \propto M^2/R^4$ is what matters, and it is robust. Two consequences to
hold onto:

\begin{enumerate}
\item \textbf{Contraction raises the required pressure steeply.} Halving $R$
  multiplies the demand by 16.
\item \textbf{The star must supply that pressure from \emph{something}.} Which
  something --- ideal gas, radiation, degenerate electrons --- is the entire
  story of stellar evolution, and in silicon burning it is degenerate electrons
  at ${\sim}$93--95\% of the total (Tier~0 \S0.1.2's worked example at
  $\rho = 10^8$, $\Tn = 4$, $\Ye = 0.5$, $\bar A = 28$; recomputed 92--95\%
  depending on the electron treatment), which is exactly why \Ye{} matters.
\end{enumerate}

% ---------------------------------------------------------------------
\subsubsection{The virial theorem, and the negative heat capacity --- derived}
\label{sec:virial}

Take the hydrostatic equation, multiply by $4\pi r^3$, and integrate over the
star (the textbook derivation --- \citealt{kippenhahn2012}, Ch.~3;
\citealt{clayton1968}, Ch.~2 --- reproduced here so the sign of the heat
capacity is not taken on trust):
%
\begin{equation}
\int_0^R 4\pi r^{3}\,\frac{dP}{dr}\,dr
\;=\; -\int_0^R \frac{G m}{r}\,4\pi r^{2}\rho\,dr \;=\; E_{\mathrm{grav}},
\label{eq:virialstart}
\end{equation}
%
where $E_{\mathrm{grav}}$ is the (negative) gravitational binding energy.
Integrating the left side by parts, with $P(R) = 0$,
%
\begin{equation}
\Big[4\pi r^{3}P\Big]_0^R - \int_0^R 12\pi r^{2} P\,dr
\;=\; -3\int_0^R \frac{P}{\rho}\,dm .
\label{eq:virialparts}
\end{equation}

So
%
\begin{equation}
\boxed{\;-3\int \frac{P}{\rho}\,dm \;=\; E_{\mathrm{grav}}.\;}
\label{eq:virial}
\end{equation}

Now specialise. For a non-relativistic ideal gas, $P/\rho = (\gamma-1)\,u$ with
$\gamma = 5/3$ and $u$ the specific internal energy, so
$\int (P/\rho)\,dm = (2/3)\,E_{\mathrm{int}}$ and
%
\begin{equation}
E_{\mathrm{int}} \;=\; -\tfrac{1}{2}E_{\mathrm{grav}},
\qquad
E_{\mathrm{tot}} \;=\; E_{\mathrm{int}} + E_{\mathrm{grav}}
\;=\; \tfrac12 E_{\mathrm{grav}} \;=\; -E_{\mathrm{int}} .
\label{eq:virialideal}
\end{equation}

\textbf{Read the last equality.} The star's total energy is \emph{minus} its
internal energy. Radiate energy away --- make $E_{\mathrm{tot}}$ more negative
--- and $E_{\mathrm{int}}$ goes \textbf{up}. The star heats when it loses
energy. Its effective heat capacity is negative.
\derivedhere[standard --- \citealt{kippenhahn2012}, Ch.~3]

This is the engine. It explains, with no further physics:

\begin{itemize}
\item why a star contracts and heats after each fuel is exhausted, with no need
  for anything to ``trigger'' the next stage;
\item why the burning stages form an ordered staircase in temperature;
\item why the process runs away at the end rather than settling;
\item and, crucially for this project, why the \emph{rate} of the sequence is
  set by the \textbf{cooling} rate. A star that loses energy faster contracts
  faster and burns faster. \S\ref{sec:neutrinoclock} turns that into the
  day-long silicon-burning timescale.
\end{itemize}

For an ultra-relativistic gas ($\gamma = 4/3$, $P/\rho = u/3$) the same algebra
gives $E_{\mathrm{int}} = -E_{\mathrm{grav}}$ and hence $E_{\mathrm{tot}} = 0$:
the star is marginally bound, and the negative heat capacity vanishes
(\citealt{shapiro1983} \S3.4 for the $\Gamma = 4/3$ marginal case).
\textbf{That is not a curiosity --- it is the collapse condition}, and
\S\ref{sec:gamma43} is the same statement with the corrections restored.

% ---------------------------------------------------------------------
\subsubsection{What temperature a contracting core reaches}
\label{sec:tcore}

Combine the virial relation with an ideal-gas equation of state,
$P = \rho k T/(\bar\mu\, m_u)$, evaluated crudely at the centre:
%
\begin{equation}
\frac{k T_c}{\bar\mu\, m_u} \;\sim\; \frac{G M}{3R}
\qquad\Longrightarrow\qquad
T_c \;\sim\; \frac{G M \bar\mu\, m_u}{3kR}
\;\propto\; M^{2/3}\rho^{1/3}.
\label{eq:tcore}
\end{equation}

The last proportionality uses $R \propto (M/\rho)^{1/3}$. Two things fall
straight out.

\textbf{The mass dependence sets which stars get anywhere.}
$T_c \propto M^{2/3}$ at fixed $\rho$: a more massive core is hotter at the same
density, which is why only stars above ${\sim}$8--10~\Msun{} reach silicon
burning at all (\citealt{heger2003}: stars below ${\sim}9$~\Msun{} end as white
dwarfs, ${\sim}$9--10~\Msun{} form degenerate ONe cores, above ${\sim}10$~\Msun{}
core collapse is the only alternative) --- below that, electron degeneracy
(\S\ref{sec:degeneracy}) halts contraction before the core is hot enough.

\textbf{The density dependence sets the staircase.} $T_c \propto \rho^{1/3}$: to
raise the core temperature by a factor of 10 you must raise its density by a
factor of 1000. Numerically, for a 1.4~\Msun{} core with $\bar\mu \approx 1.87$
(fully ionised \nuc{Si}{28}: $\mu = A/(Z+1) = 28/15$; an earlier version used
$\bar\mu \approx 1.4$, which is no silicon-ash composition --- corrected in the
2026-08-18 audit): \derivedhere

\begin{center}
\begin{tabular}{@{}ll@{}}
\toprule
\textbf{$\rho_c$ [\gcc]} & \textbf{$T_c$ from the virial estimate [GK]} \\
\midrule
$10^{2}$ & 0.074 \\
$10^{6}$ & 1.59 \\
$10^{8}$ & 7.4 \\
\bottomrule
\end{tabular}
\end{center}

The last row is the punchline: \textbf{a Chandrasekhar-mass core compressed to
${\sim}10^8$~\gcc{} is, by this estimate, at 7.4~GK --- silicon-burning
temperature to within the factor ${\sim}2$ that the ideal-gas assumption
costs.} The actual ${\sim}3.5$~GK is the regime box's centre; the crude estimate
lands at its upper edge. The pairing was not chosen --- it is where the star
puts itself. The box (\Tn{} 1.6--7.9, $\rho$ $10^7$--$10^9$, Tier~0 \S III.2)
spans $\pm 1$ decade in $\rho$ and $\pm 0.35$~dex in $T$ around
($10^8$~\gcc, 3.5~GK).

The estimate is crude by a factor of ${\sim}2$ in $T$ (7.4~GK against the
${\sim}3.5$~GK at which silicon actually burns; the ideal-gas assumption is
already wrong at $10^8$~\gcc, where electrons are degenerate ---
\S\ref{sec:degeneracy}), and it should not be quoted as a prediction. It is
here to show that the regime box is \emph{derivable} from $M$, $G$ and $k$, not
sampled from a paper.

% ---------------------------------------------------------------------
\subsubsection{Degeneracy: the wall the star runs into}
\label{sec:degeneracy}

Electron degeneracy pressure is independent of temperature (in the $T \to 0$
limit; \citealt{kippenhahn2012}, Ch.~15; \citealt{shapiro1983} \S2.3), so a
degenerate core does \textbf{not} heat when it contracts --- the virial argument
above assumed an ideal gas and fails. Whether a core is degenerate is decided by
comparing the electron Fermi energy with $kT$ (the zero-temperature Fermi
momentum and kinetic Fermi energy in the form of \citealt{chabrier1998}
\S II):
%
\begin{equation}
\eF \;=\; \sqrt{(p_F c)^2 + (m_e c^2)^2} - m_e c^2,
\qquad
p_F = \hbar\left(3\pi^{2} n_e\right)^{1/3},
\qquad
n_e = \rho N_A Y_e .
\label{eq:fermienergy}
\end{equation}

At $\Tn = 4$ ($kT = 0.345$~MeV), $\Ye = 0.5$ ($\Ye = 0.46$ lowers \eF{} by
3--5\%): \derivedhere

\begin{center}
\begin{tabular}{@{}llll@{}}
\toprule
\textbf{$\rho$ [\gcc]} & \textbf{\eF{} [MeV]} & \textbf{$\eF / kT$} &
  \textbf{regime} \\
\midrule
$10^{6}$ & 0.14 & 0.4  & non-degenerate \\
$10^{7}$ & 0.51 & 1.5  & marginal \\
$10^{8}$ & 1.46 & 4.2  & degenerate \\
$10^{9}$ & 3.61 & 10.5 & strongly degenerate \\
\bottomrule
\end{tabular}
\end{center}

So the regime box spans the degeneracy transition --- its low-density edge is
only marginally degenerate ($\eF/kT \approx 1.5$) and its high-density edge is
strongly so. That single fact is why the box is a box rather than a line, why
the equation of state has four competing components (Tier~0 \S0.1.1), and why
the electron-capture threshold (which needs \eF{} above a nuclear threshold) is
a \emph{box edge} rather than a global on/off switch (Tier~0 \S0.1.5).

\textbf{The astrophysical role of degeneracy.} A star that becomes degenerate
before igniting its next fuel is stuck: contraction no longer heats it, so it
cools into a white dwarf. That is the fate of everything below
${\sim}$8--10~\Msun{} \citep{heger2003}. Above that mass the core stays hot
enough, relative to \eF, to keep igniting --- until the iron core, where there
is no next fuel and degeneracy is all that is holding the star up. \textbf{The
whole massive-star story is a race between contraction-heating and degeneracy,
and the iron core is where degeneracy finally wins and then immediately loses}
(\S\ref{sec:corecollapse}).

% =====================================================================
\subsection{The burning staircase, derived rather than tabulated}
\label{sec:staircase}

Tier~0 \S I.1 tabulates the stages. This section derives \emph{why the table
looks like that}, because three of the project's design facts --- the
Si-burning temperature, the one-day duration, and the neutrino-dominated
energetics --- are consequences of the derivation rather than of the table.

% ---------------------------------------------------------------------
\subsubsection{Why each fuel ignites where it does}
\label{sec:ignition}

Tier~1 \S1.2 derives the Gamow peak (textbook: \citealt{clayton1968} \S4-3;
\citealt{iliadis2015} \S3.2.1; the $E_0/kT$ and $\exp(-3E_0/kT)$ forms are also
in \citealt{adelberger2011} \S II.A): for a non-resonant charged-particle
reaction between nuclei of charges $Z_1$, $Z_2$ and reduced mass $\mu$, the
effective burning energy and the rate's temperature dependence are governed by
%
\begin{equation}
E_0 \;=\; \left(\frac{b\,kT}{2}\right)^{2/3},
\qquad
b \;=\; \frac{\sqrt{2\mu}\,\pi Z_1 Z_2 e^{2}}{\hbar},
\qquad
\langle\sigma v\rangle \;\propto\;
\exp\!\left[-\,3\left(\frac{b}{2}\right)^{2/3}\!(kT)^{-1/3}\right].
\label{eq:gamowpeak}
\end{equation}

It is convenient to name the combination $b^2 \equiv \EG$, the \textbf{Gamow
energy}, and to write the exponent as a single dimensionless number:
%
\begin{equation}
\EG \;=\; 2\mu c^{2}\left(\pi\alpha Z_1Z_2\right)^{2}
\;=\; 0.979\,\left(Z_1Z_2\right)^{2}A_{\mathrm{red}}\ \mathrm{MeV},
\qquad
\tau \;\equiv\; \frac{3E_0}{kT} \;=\; 3\left(\frac{\EG}{4kT}\right)^{1/3},
\label{eq:gamowenergy}
\end{equation}
%
so that $\langle\sigma v\rangle \propto e^{-\tau}$ (the coefficient
$0.979\,(Z_1Z_2)^2 A_{\mathrm{red}}$~MeV is \citet{rolfs1988}'s; CODATA
$\alpha$ gives 0.9791). Ignition happens when the rate becomes fast enough to
balance the star's losses. \textbf{The tempting move at this point is wrong and
worth doing carefully}, because getting it wrong is how one ends up with a
staircase that does not exist.

The tempting move: hold the exponent fixed at ignition,
$\tau_{\mathrm{ign}} \approx \mathrm{const}$, and read off
$\Tign \propto \EG \propto (Z_1Z_2)^2\mu$. That would make carbon ignite at
$36^2 \times 12 \approx 1.6 \times 10^4$ times the pp ignition temperature, i.e.
at ${\sim}60$~GK --- hotter than anything short of the bounce shock. It does
not; the actual factor is 200. Inverting the definition of $\tau$ instead,
%
\begin{equation}
k\Tign \;=\; \frac{\EG}{4}\left(\frac{3}{\tau_{\mathrm{ign}}}\right)^{3}
\qquad\Longrightarrow\qquad
\Tign \;\propto\; \frac{\EG}{\tau_{\mathrm{ign}}^{3}} ,
\label{eq:tign}
\end{equation}
%
and \textbf{$\tau_{\mathrm{ign}}$ is not slowly varying} --- it is the whole
story. Check against the fuels, with \EG{} computed from the formula above and
$\tau_{\mathrm{ign}}$ evaluated at each stage's actual ignition temperature:
\derivedhere[ignition temperatures after \citealt{woosley2002} (not re-read
here), bracketed by \citealt{kato2020} for a 15~\Msun{} MESA model --- H
$4 \times 10^7$, He $1.5 \times 10^8$, C $7 \times 10^8$, Ne $1.4 \times 10^9$,
O $1.6 \times 10^9$, Si ${\sim}3 \times 10^9$~K --- and by \citealt{arnett1989}
for 20~\Msun{} (via \citealt{odrzywolek2004}): C 0.81, Ne 1.69, O 2.10,
Si 3.70~GK]

\begingroup\small
\begin{tabularx}{\textwidth}{@{}L{1.3cm} L{2.2cm} L{0.9cm} L{0.9cm} L{1.5cm} Y L{1.2cm}@{}}
\toprule
\textbf{fuel} & \textbf{dominant channel} & \textbf{$Z_1Z_2$} &
  \textbf{$A_{\mathrm{red}}$} & \textbf{\EG{} [MeV]} &
  \textbf{\Tign{} (actual)} & \textbf{$\tau_{\mathrm{ign}}$} \\
\midrule
H (pp) & p + p & 1 & 0.50 & 0.49 &
  0.004~GK (pp threshold in the lowest-mass stars; H ignites at
  0.03--0.04~GK in a massive star, \citealt{kato2020}) & 21.3 \\
H (CNO) & \nuc{N}{14}(p,$\gamma$)\nuc{O}{15} & 7 & 0.933 & 44.8 & 0.03~GK &
  48.9 \\
He & $\alpha + \alpha$ ($\to$ \nuc{C}{12}) & 4 & 2.0 & 31.3 & 0.15--0.2~GK &
  \emph{resonant} \\
C & \nuc{C}{12} + \nuc{C}{12} & 36 & 6.0 & 7\,617 & 0.8~GK & 90.7 \\
Ne & \nuc{Ne}{20}($\gamma$,$\alpha$) & --- & --- & --- (\textbf{photo}) &
  1.5~GK & --- \\
O & \nuc{O}{16} + \nuc{O}{16} & 64 & 8.0 & 32\,100 & 2.0~GK & 107.9 \\
Si & \nuc{Si}{28}($\gamma$,$\alpha$) & --- & --- & --- (\textbf{photo}) &
  3--4~GK & --- \\
\bottomrule
\end{tabularx}
\endgroup

Because $\tau_{\mathrm{ign}}$ is \emph{defined} from the actual \Tign,
$\EG/\tau_{\mathrm{ign}}^3 = 4k\Tign/27$ reproduces the ignition temperatures
\emph{identically} --- that is a bookkeeping identity, not a test (an earlier
version of this paragraph presented the agreement as a prediction ``to better
than 1\%''; corrected in the 2026-08-18 audit). What the table does show is
\emph{how the ratio is made}: from pp to oxygen, \textbf{\EG{} grows by
$6.6 \times 10^4$ while \Tign{} grows by only 500, so $\tau_{\mathrm{ign}}^3$
must absorb a factor of 130} ($\tau_{\mathrm{ign}}$ rising from 21 to 108). The
Coulomb barrier really is enormous; the reason the staircase is compressed into
under three decades of temperature rather than nearly five is that the star
simply burns deeper into the exponential tail for each successive fuel
($\tau_{\mathrm{ign}} \approx 21 \to 49 \to 91 \to 108$). What sets
$\tau_{\mathrm{ign}}$ is the rate the star needs --- $E_{\mathrm{fuel}}$ divided
by the loss rate --- which is why it is not universal.

Two rows do not belong to that scaling at all, for two different reasons.

\textbf{Helium is resonant.} Triple-$\alpha$ proceeds through the \nuc{Be}{8}
ground state and the Hoyle resonance in \nuc{C}{12} (\citealt{salpeter1952};
\citealt{hoyle1954}; textbook, \citealt{iliadis2015} \S5.2.2), so its rate is
set by a narrow-resonance formula, not by the non-resonant Gamow integral; its
$\tau$ is not the controlling quantity and the row is listed only for the
barrier comparison.

\textbf{Neon and silicon are not fusion at all.} They do not ignite because
their Coulomb barriers become penetrable; they ignite because the thermal photon
bath acquires enough $\gamma$-rays above the ($\gamma$,$\alpha$) threshold
(\citealt{kato2020} for neon; \citealt{boccioli2024} for the ``re-adjustment of
the chemical composition'' character of the O-to-Fe stages). Take the same
estimate seriously for \nuc{Si}{28} + \nuc{Si}{28}: $Z_1Z_2 = 196$ and
$A_{\mathrm{red}} = 14$ give $\EG = 5.3 \times 10^5$~MeV, $16.4\times$
oxygen's, and extrapolating the slowly-rising
$\tau_{\mathrm{ign}} \approx 115$--130 puts fusion ignition at
\textbf{${\approx}20$~GK} (19--27~GK across that $\tau$ range) --- far above
where the star ever gets in hydrostatic burning, and above the temperature at
which the Fe peak itself dismantles. More directly: at $\Tn = 3.5$ the Si + Si
exponent is $\tau = 228$ against 108 for O + O at its ignition, i.e. a
penetrability suppressed by $e^{-120}$.

\textbf{This is the single most important structural fact in the project}, and
it is worth stating in the sharpest available form:

\begin{invariant}
\nuc{Si}{28} cannot burn by fusing with itself, at any temperature the star
survives. It burns by being taken apart.
\end{invariant}

Everything downstream --- the near-balanced forward/reverse pairs, the
quasi-equilibrium clustering, the catastrophic cancellation, the $\kappa$
diagnostic, the stiffness, and the entire Target A design --- is a consequence
of that sentence.

The photodisintegration threshold itself is derivable. The relevant quantity is
$Q/kT$ for the ($\gamma$,$\alpha$) channel, with $kT = 86.2~\mathrm{keV} \times
\Tn$:
%
\begin{equation}
{}^{28}\mathrm{Si}(\gamma,\alpha)^{24}\mathrm{Mg}: \quad Q = 9.98\ \mathrm{MeV}
\qquad\Longrightarrow\qquad
\frac{Q}{kT} = \frac{115.8}{T_9}.
\label{eq:photothreshold}
\end{equation}

At $\Tn = 3$ that exponent is 38.6; at $\Tn = 4$ it is 28.9.
\derivedhere[Q from the AME2020 mass excesses
($S_\alpha(\nuc{Si}{28}) = 9.984$~MeV), Tier~0 \S0.3.1] Between those two
temperatures the Boltzmann suppression changes by a factor
$e^{9.65} \approx 1.6 \times 10^4$. The transition from ``no
photodisintegration'' to ``photodisintegration controls everything'' occupies
less than a factor of 1.5 in temperature, which is why \S0.2 of Tier~2 can say
the regime change happens across less than a factor of two in $T$, and why the
kill-test priority window is 3.3--5~GK.

% ---------------------------------------------------------------------
\subsubsection{The energy per gram of each stage, computed}
\label{sec:energypergram}

Take the binding energies (Tier~0 \S0.3.1's sign convention: $B > 0$, and the
energy released is $\Delta B$) and compute the specific energy release of each
stage. For a reaction with total $\Delta B$ over $A$ nucleons, the release per
gram is $N_A \times (\Delta B/A)$~[MeV] $\times\ 1.602 \times 10^{-6}$~erg/MeV.
\derivedhere[Q-values from the AME2020 mass excesses \citep{wang2021} with
CODATA 2018 constants \citep{tiesinga2021}, all four rows recomputed in the
2026-08-18 audit]

\begingroup\small
\begin{tabularx}{\textwidth}{@{}L{1.7cm} Y L{2.6cm} L{2.0cm} L{2.2cm}@{}}
\toprule
\textbf{stage} & \textbf{representative reaction} & \textbf{$\Delta B$ [MeV]} &
  \textbf{per nucleon [MeV]} & \textbf{per gram [erg/g]} \\
\midrule
H $\to$ He & $4p \to \nuc{He}{4} + 2e^+ + 2\nu$ &
  26.7 (gross; $\nu$ carry off ${\sim}$2--7\%) & 6.7 & $6.5 \times 10^{18}$ \\
C burning & $2\,\nuc{C}{12} \to \nuc{Ne}{20} + \alpha$ & 4.62 & 0.192 &
  $1.9 \times 10^{17}$ \\
O burning & $2\,\nuc{O}{16} \to \nuc{Si}{28} + \alpha$ & 9.60 & 0.300 &
  $2.9 \times 10^{17}$ \\
Si burning & $2\,\nuc{Si}{28} \to \nuc{Ni}{56}$ & 10.92 & 0.195 &
  $1.9 \times 10^{17}$ \\
\bottomrule
\end{tabularx}
\endgroup

Two lessons.

\textbf{(i) Hydrogen burning releases 22--35$\times$ more per gram than anything
after it} (34$\times$ against carbon and silicon, 22$\times$ against oxygen).
The first stage carries almost the whole nuclear energy budget of the star's
life; everything from carbon onward is scraping the barrel. This is why the
late stages cannot last: there is very little fuel value left, and
(\S\ref{sec:neutrinoclock}) the losses are enormous.

\textbf{(ii) Silicon burning is not even the most productive late stage.} Oxygen
burning releases 50\% more per gram. Silicon's distinction is not its energy
yield --- it is that it is the \emph{last} one, that it produces the iron peak,
and that it is where the composition acquires the neutron excess that decides
what happens next.

For the record, the binding-energy-per-nucleon values that matter at the top of
the curve: \nuc{Fe}{56} 8.790, \nuc{Ni}{62} 8.795, \nuc{Ni}{56}
8.643~MeV/nucleon. \provtag{[AME2020: 8790.36, 8794.56, 8642.78 keV]}
\nuc{Ni}{62} is the maximum, not \nuc{Fe}{56} \citep{fewell1995} --- the
standard shorthand is false, and Tier~2 \S0.2.3 explains why \nuc{Fe}{56}
nonetheless wins in NSE at $\Ye \approx 0.46$: the free energy is minimised at
fixed \Ye, and the winner is the nucleus whose $Z/A$ \emph{matches} \Ye{} among
the near-maximally-bound ones.

% ---------------------------------------------------------------------
\subsubsection{The neutrino clock --- why the last stages are short}
\label{sec:neutrinoclock}

Here is the fact that turns an ordered staircase into an accelerating collapse.

\textbf{From carbon burning onward, the star's dominant energy loss is neutrino
emission from the plasma, not photon radiation from the surface}
(\citealt{limongi2017}: pair-neutrino emission ``starts to become efficient when
the central temperature exceeds ${\sim}8 \times 10^8$~K, i.e., at the beginning
of core C burning''; \citealt{odrzywolek2004}). The contrast is in \emph{how the
energy gets out}, not in how long a fuel lasts. Photons are trapped: the
interior is optically thick, so energy random-walks outward on a thermal
(Kelvin--Helmholtz) timescale of ${\sim}10^4$ years for a massive main-sequence
star (much shorter for the extended supergiant envelope), is reprocessed on the
way, and emerges from the surface at whatever rate the star's structure permits
--- the luminosity is set by the envelope, not by the burning region. Neutrinos
have cross-sections around $10^{-44}$~cm$^2$ at these energies
(\citealt{odrzywolek2004}: spectrum-averaged
$\sigma(\bar\nu_e p) \approx 7 \times 10^{-44}$~cm$^2$ for pre-supernova
neutrinos) and mean free paths vastly exceeding the stellar radius, so they leave
\textbf{at the speed of light, directly from the point of production} (Tier~0
\S0.4.2; \citealt{limongi2017}), and the loss rate is set by the local plasma
conditions alone. The core is optically thin to its own dominant coolant, and
that is what decouples its cooling rate from its structure.

(The \emph{duration} of each stage is $E_{\mathrm{fuel}}/L$, a nuclear timescale
--- $10^7$ years for hydrogen (\citealt{kato2020} for a 15~\Msun{} model;
\citealt{limongi2017}: $10^7$--$10^6$~yr over 13--120~\Msun) because there is a
great deal of hydrogen and \S\ref{sec:energypergram}'s release per gram is
large, not because of any transport time. What the transport argument
establishes is the switch of coolant, and \S\ref{sec:neutrinoclock}'s remaining
paragraphs are what turn that into a duration.)

The dominant thermal channel here is \textbf{electron--positron pair
annihilation}, $e^+e^- \to \nu\bar\nu$ (\citealt{kato2020};
\citealt{patton2017}). Its emissivity requires a thermal positron population,
which requires $kT$ no longer negligible against $m_e c^2 = 0.511$~MeV ---
$kT/m_e c^2 \approx 0.17$ at $\Tn = 1$, where the Boltzmann factor
$e^{-2m_e c^2/kT} \approx 10^{-5}$ already makes the pair population
non-negligible; onset ${\sim}$0.8--1~GK (\citealt{limongi2017};
\citealt{odrzywolek2004}). Once that threshold is crossed the positron number
density rises as roughly $\exp(-2m_e c^2/kT)$ times phase-space factors, and
the emissivity per gram rises extremely steeply --- approximately as $T^9$ in
the relativistic non-degenerate limit (\citealt{fowler1964};
\citealt{itoh1996}, ApJS 102, 411, for the fitting formulae stellar codes use);
across 2--3.5~GK the effective exponent of the pair rate is ${\approx}$7--8
\derivedhere[from the non-relativistic form
$\varepsilon \propto \Tn^3\, e^{-11.86/\Tn}$], and the emissivity per gram is
further reduced by electron degeneracy at $\rho \gtrsim 10^7$--$10^8$~\gcc{}
(\citealt{kato2020} \S2).

Put the two facts together with the virial engine of \S\ref{sec:virial}:
%
\begin{equation}
\begin{gathered}
\text{contraction}
\;\longrightarrow\; T\uparrow
\;\longrightarrow\; \varepsilon_\nu \propto T^{9}\ \text{rises catastrophically}\\[2pt]
\;\longrightarrow\; \text{must burn faster}
\;\longrightarrow\; \text{contraction}
\end{gathered}
\label{eq:coolingrunaway}
\end{equation}

The runaway is not in the burning; it is in the \textbf{cooling}
(\citealt{limongi2017}: ``an almost constant nuclear energy \ldots{} coupled to
the dramatic increase of the total luminosity''). Each stage must generate
energy at the rate the neutrinos remove it, and that rate rises by roughly seven
to nine powers of $T$ while the available fuel energy per gram
(\S\ref{sec:energypergram}) is roughly constant. So the duration of stage $n$
scales as
%
\begin{equation}
\tau_n \;\sim\; \frac{E_{\mathrm{fuel}}}{\varepsilon_\nu(T_n)} \;\propto\; T_n^{-9}
\label{eq:stageduration}
\end{equation}
%
to leading order (effective exponent closer to 7--8 across 2--3.5~GK; $\rho$
and degeneracy neglected). Check it: from oxygen burning ($\Tn \approx 2$,
months --- 180~d for 20~\Msun, \citealt{arnett1989} via
\citealt{odrzywolek2004}; 14 months for 15~\Msun, \citealt{kato2020}) to silicon
burning ($\Tn \approx 3.5$), the temperature ratio is 1.75 and $1.75^9 = 154$
(with the effective exponent 7.5 the factor is $1.75^{7.5} \approx 70$ --- still
``of order a day''). Months\,/\,154 is of order a day. \textbf{The one-day
silicon-burning timescale is a $T^9$ scaling law applied to the oxygen-burning
duration.} \derivedhere[{order-of-magnitude only; the actual duration depends
on core mass and shell structure: roughly a day (shorter for the most massive
cores) to a couple of weeks --- 2~d (20~\Msun; \citealt{arnett1989} via
\citealt{odrzywolek2004}), 4.5~d (15~\Msun{} MESA model; \citealt{kato2020}),
${\sim}10^{-2}$~yr ${\approx}$ 4~d \citep{limongi2017}; the ${\sim}$2-week upper
end is Tier~0 \S I.1 (after \citealt{woosley2002}), not re-verified here
\typical}]

Three consequences that shape this project:

\begin{enumerate}
\item \textbf{The regime is transient and the star never settles.} There is no
  steady-state silicon-burning configuration to compute; the whole episode is a
  relaxation with continuously changing $(T, \rho)$. That is why the training
  data is a grid of \emph{initial states} burned for fixed $\Delta t$ rather
  than a library of equilibria, and why the operator-splitting framing of
  Tier~0 \S IV.1 is the right one.
\item \textbf{Neutrino losses are a first-class output, not an afterthought.}
  \epsnu{} appears in the NNN's loss function and in the emulator's specified
  output head for the same reason it appears in the star's energy equation: it
  is the dominant term. Tier~1 \S VIII.C.3 flags that the flux route currently
  drops the tabulated NU column; \S\ref{sec:neutrinoclock} is why that flag
  matters.
\item \textbf{``Burning'' and ``cooling'' are not separable.} The $T^9$ coolant
  and the temperature-sensitive burning are locked together. An emulator
  validated at \emph{imposed} $(T, \rho)$ --- which is all of this project's
  data --- has never been tested inside that loop. This is
  \tierone{VIII.C.4}\,/\,\foreignreg{2}{R-11}, and \S\ref{sec:mesastep}
  returns to it.
\end{enumerate}

% ---------------------------------------------------------------------
\subsubsection{What the star is doing while the network is being called}
\label{sec:whilecalled}

To ground the above: at silicon-burning conditions, the star's core is
${\sim}1.4$~\Msun{} of matter a few thousand km in radius, with central density
$\rho_c \sim 10^8$~\gcc{} ($4.9 \times 10^7$ in the 20~\Msun{} model of
\citealt{arnett1989} via \citealt{odrzywolek2004}; a uniform 1.4~\Msun{} sphere
at $10^8$ would be only $1.9 \times 10^3$~km, a centrally condensed core with
$\rho_c \sim 10^8$ is ${\sim}7 \times 10^3$~km --- an earlier version wrote ``a
${\sim}10^4$~km sphere at $\rho \sim 10^8$'', an inconsistent triple; corrected
in the 2026-08-18 audit) and $T_c \sim 3.5$~GK (3.7~GK in that model),
converting silicon to iron-peak nuclei over about a day
(\citealt{farmer2016}: ``Silicon ignites at the center within one day of
core-collapse''), radiating almost all of that energy as neutrinos that leave
instantly, contracting as the fuel is depleted and the entropy is drained by
neutrinos, and building an iron core that grows toward a mass limit it cannot
exceed.

Meanwhile, in a MESA calculation of the same object, the nuclear network is
being evaluated for every zone at every Newton iteration of every timestep in
the default fully coupled mode (once per hot zone per step, with internal
sub-steps, in \code{op\_split\_burn} mode) --- hundreds to thousands of zones
(\citealt{paxton2011} \S6.1), and timesteps so short at the end that the last
day of the star's life can consume a large fraction of the total run time of the
whole $10^7$-year calculation. \S\ref{part:host} quantifies that.

% ---------------------------------------------------------------------
\subsubsection{Convection --- the transport a one-zone burner cannot see}
\label{sec:convection}

Tier~0 \S0.7 puts the convective turnover time in the timescale table
($10^1$--$10^3$~s, with an honest warning that both factors in
$L/v_{\mathrm{conv}}$ are uncertain by a decade; the 3D Si- and O-shell
simulations of \citealt{couch2015} and \citealt{muller2016b} give
${\approx}20$~s) and concludes ``one-zone is an approximation''. That is where
the treatment stops in every existing tier. It should not, because mixing is
not a small correction to late-stage burning --- \textbf{it is co-dominant with
it}, and it changes what the emulator's error budget means.

\textbf{Silicon shell burning is convective} (\citealt{couch2015}: ``strong
convection driven by violent Si burning in the shell surrounding the iron
core''). The energy generation rate is a steep function of temperature, so the
burning region is strongly superadiabatic and convection sets in immediately. In
a 1D stellar-evolution code this is modelled by mixing-length theory (MLT):
convection becomes a diffusive process with diffusivity
$D \approx \tfrac13\, v_{\mathrm{conv}}\,\ell$ (\citealt{paxton2011} \S5.2;
MESA r23.05.1 \code{turb/private/mlt.f90:156},
\code{D = conv\_vel*Lambda/3d0}), and the composition equation acquires a
transport term (\citealt{paxton2011} \S6):
%
\begin{equation}
\frac{\partial Y_i}{\partial t}
\;=\; \underbrace{\big[\nu\,\mathbf R(\mathbf Y;T,\rho)\big]_i}_{\text{the network --- what this project emulates}}
\;+\;\underbrace{\frac{\partial}{\partial m}\!\left[(4\pi r^2\rho)^2 D\,\frac{\partial Y_i}{\partial m}\right]}_{\text{mixing --- outside the split}} .
\label{eq:burnmix}
\end{equation}

Those two terms can be \textbf{operator-split} from each other. With MESA's
\code{op\_split\_burn} option (\citealt{jermyn2023} \S10.2 --- off by default,
but switched on in MESA's own core-collapse test-suite models and the mode the
NNN paper \S4 envisages plugging into; the default solver couples burning,
mixing and structure in one Newton solve, \citealt{paxton2011} \S6.1), each hot
zone is burned first at fixed $(T, \rho)$, and mixing is then solved together
with the structure. The emulator replaces the burn operator only. (An earlier
version described the burn-mix-structure sequence as MESA's universal behaviour;
corrected in the 2026-08-18 audit.)

\textbf{The competition is quantified by a Damk\"ohler number.} Define
%
\begin{equation}
\Da_i \;=\; \frac{\tau_{\mathrm{mix}}}{\tau_{\mathrm{burn},i}} ,
\label{eq:damkohler}
\end{equation}
%
the ratio of the mixing time to the timescale on which process $i$ changes the
composition (the standard combustion definition,
$\Da = \tau_{\mathrm{transport}}/\tau_{\mathrm{reaction}}$, so $\Da > 1$ means
burning is faster than mixing). $\Da \gg 1$ means burning wins and gradients
survive; $\Da \ll 1$ means mixing wins and the region is homogeneous in that
quantity. Read the numbers off Tier~0 \S0.7's table with
$\tau_{\mathrm{mix}} \approx 10^1$--$10^3$~s: \derivedhere

\begin{center}
\begin{tabularx}{\textwidth}{@{}L{4.2cm} L{2.6cm} L{2.8cm} Y@{}}
\toprule
\textbf{process} & \textbf{$\tau$} & \textbf{$\Da = \tau_{\mathrm{mix}}/\tau$} &
  \textbf{verdict} \\
\midrule
intra-group equilibration & $\lesssim 10^{-3}$~s &
  \textbf{$10^{4}$ -- $10^{6}$} & burning wins overwhelmingly \\
inter-group (bridge) flow & $10^{-1}$ -- $10^{2}$~s &
  \textbf{$10^{-1}$ -- $10^{4}$} & \emph{comparable} --- no clean answer \\
weak reactions\,/\,\Ye & $10^{2}$ -- $10^{5}$~s &
  \textbf{$10^{-4}$ -- $10^{1}$} & \textbf{mixing wins, mostly} \\
\bottomrule
\end{tabularx}
\end{center}

Three consequences, and the second is the one that matters most.

\textbf{(i) The fast species are local; \Ye{} is not.} Group-internal
equilibrium is established far faster than material can be transported, so the
QSE composition is a \emph{local} function of the local state --- which is what
licenses treating the network call as a per-zone problem at all. But \Ye{}
evolves on a timescale comparable to or longer than mixing, so \textbf{\Ye{} is
spatially homogenised across the convective region}. The quantity the whole
project is organised around is the one mixing most strongly averages.

\textbf{(ii) That averaging does not save a systematic error --- and the
arithmetic is \S\ref{sec:mch}'s argument, in space instead of time.} Suppose the
emulator makes a per-zone \Ye{} error $\delta$, and the convective region spans
$N_z$ zones (of order $10^2$--$10^3$ \assumednote{order of magnitude;
\citet[\S 6.1]{paxton2011} give ``hundreds to thousands'' of cells for the
\emph{whole} model} in a MESA model). After mixing, the region's error is
%
\begin{equation}
\delta_{\text{mixed}} \;=\;
\begin{cases}
\delta/\sqrt{N_z} & \text{if per-zone errors are independent},\\[2pt]
\delta & \text{if the error is a systematic function of the state.}
\end{cases}
\label{eq:mixederror}
\end{equation}

Zones in one convective region have \emph{similar} $(T, \rho, X)$ --- that is
what being in one convective region means --- so a learned model's error on
them is strongly correlated, not independent. \textbf{Mixing averages away
almost none of a neural network's error, while it would have averaged away most
of a random one.} A third independent argument for the conclusion Tier~2
\S0.4.2 reaches for soft penalties and Tier~0 \S IV.3 reaches for rollout:
\emph{the character of the error matters more than its size, and correlated
error is the enemy.} \derivedhere

\textbf{(iii) Mixing feeds the network states that are mixtures, not
trajectories.} The composition arriving at a network call is a blend of burnt
and unburnt material from different depths --- not the endpoint of any single
burning history. This cuts \emph{in favour} of the Sobol design in a way nobody
in the project has said out loud:

\begin{warn}
The relaxed manifold (\S\ref{sec:twodistributions}, distribution B) is built
from constant-$(T, \rho)$ single-history relaxations, the \emph{opposite}
extreme from the training grid's random mixtures. Real convective burning sits
\textbf{between} them: compositions partially relaxed \emph{and} partially
mixed. Neither measured distribution is the deployment distribution, and the
true one is probably not bracketed by them either, because mixing can inject
fuel no relaxation would produce.
\end{warn}

That is a sharper statement of the distribution problem than
\S\ref{sec:twodistributions} gives alone, and it means the Sobol $\to$
real-MESA gate is testing something more complicated than ``physical versus
unphysical states''.

\textbf{And the modelling caveat that tempers all of it.} MLT captures the bulk
velocity scale of late-stage shell convection reasonably well
(\citealt{meakin2007}; \citealt{muller2016b}) but misses the boundary physics.
Three-dimensional hydrodynamic simulations of oxygen and silicon shell burning
show turbulent entrainment at convective boundaries (\citealt{meakin2007};
\citealt{muller2016b}: the O shell grew from 0.51 to 0.56~\Msun{} by
entrainment), large-scale bulk motions at Mach numbers approaching 0.1 in the
oxygen shell at collapse (\citealt{muller2016b}; ${\sim}0.03$ in
\citealt{meakin2007}; ${\sim}0.02$ in the Si shell of \citealt{couch2015} as
estimated by \citealt{muller2016b}), and --- in some models --- \textbf{shell
mergers}, where the oxygen and neon shells convectively merge and the
composition changes qualitatively (in 1D: \citealt{sukhbold2018}; 3D
``encroachment'' short of a full merger before collapse: \citealt{muller2016a};
the full 3D O--Ne merger of \citealt{yadav2020} and the 1D shell-merger survey
of \citealt{collins2018} were not opened here). Shell mergers change the
pre-supernova structure and the compactness (\citealt{sukhbold2018}: ``when this
occurs, the compactness is usually small'') and the yields
\provtag{\textsc{[assumed]}}, and 1D codes reproduce them only through
calibrated prescriptions. (An earlier version said MLT ``is known to describe
late-stage convection poorly''; both cited 3D papers say the opposite of the
bulk flow --- corrected in the 2026-08-18 audit.)

So the honest hierarchy of modelling error in a pre-supernova model is
\textbf{3D convection $\gg$ network size $\gtrsim$ rate uncertainties $\gg$
emulator target} (the network-size rung: \citealt{farmer2016} find the isotope
count ``plays a more significant role in determining the span of the variations
for neon, oxygen and silicon burning'' than resolution; the rate-uncertainty
rung is \citealt{fields2018}, not opened here; the ordering itself is this
author's judgement) --- \S\ref{sec:errorhierarchy} assembles it properly,
because it determines what the project can claim.

Nothing about mixing appears in any tier's register. Registered as \reg{T-19},
and it feeds \reg{T-20} (the error hierarchy).

% =====================================================================
\subsection{Self-check}
\label{sec:selfcheck-engine}

Answer without scrolling up.

\begin{selfcheck}
% source 0.10 Q1
\item Derive the virial statement $E_{\mathrm{tot}} = -E_{\mathrm{int}}$ for a
  non-relativistic ideal-gas star, and say in one sentence why it makes stellar
  evolution a \emph{sequence} rather than an equilibrium.
% source 0.10 Q2
\item $T_c \propto M^{2/3}\rho^{1/3}$. Use it to explain both why 8~\Msun{} is
  roughly the threshold for reaching silicon burning and why the regime box sits
  at $\rho \sim 10^8$.
% source 0.10 Q3
\item Why does \nuc{Si}{28} not burn by fusing with \nuc{Si}{28}, and what
  single structural consequence does that have for every equation in
  Tiers~1--2?
% source 0.10 Q4
\item Silicon burning releases \emph{less} energy per gram than oxygen burning.
  Why is it nevertheless the stage this project emulates?
% source 0.10 Q5
\item The last stages of a massive star's life accelerate. Name the coolant, its
  temperature scaling, and use it to get the silicon-burning duration from the
  oxygen-burning duration.
% source 0.10 Q11
\item Compute the Damk\"ohler number for intra-group equilibration, bridge flow
  and the weak sector, and say which composition quantity convection
  homogenises.
% source 0.10 Q12
\item A convective region spans 300 zones. By how much does mixing reduce the
  emulator's \Ye{} error if it is random? If it is a smooth function of the
  state? What does that add to the case for structural conservation?
\end{selfcheck}
